\documentclass[11pt]{article}
\usepackage[margin=1in]{geometry}\usepackage{amsmath,amssymb,amsthm,booktabs,url}
\newtheorem{theorem}{Theorem}\newtheorem{lemma}[theorem]{Lemma}\newtheorem{proposition}[theorem]{Proposition}
\theoremstyle{remark}\newtheorem{remark}[theorem]{Remark}
\title{Counterexamples to two conjectures on domination and minimum maximal matchings in regular graphs}
\author{Andrew Butz\thanks{andrew@cosmicbutz.com. The author gratefully acknowledges extensive AI assistance (Anthropic's Claude) in computation, proof development and verification.}}
\date{September 2026}
\begin{document}\maketitle
\begin{abstract}
Baste, F\"urst, Henning, Mohr and Rautenbach~\cite{BFHMR} conjectured that every regular graph $G$ of positive degree satisfies $\gamma(G)\le\gamma_e(G)$, where $\gamma$ is the domination number and $\gamma_e$ the edge domination number, the minimum size of a maximal matching; the conjecture is known for degree at least seven~\cite{Gupta} and for every cubic graph on at most $48$ vertices~\cite{Gupta}. The stronger statement $i(G)\le\gamma_e(G)$, with $i$ the independent domination number, was proposed by the TxGraffiti system in 2020 and is Conjecture~3 of~\cite{DBP}. We exhibit six pairwise non-isomorphic connected cubic graphs on $50$ vertices with $i(G)=16>15=\gamma_e(G)$, three of which also have $\gamma(G)=16$. Thus both conjectures fail for cubic graphs, and by the $48$-vertex bound of~\cite{Gupta} the three graphs are counterexamples of minimum order to the first. Lexicographic products with edgeless graphs turn any of the six into a connected $3k$-regular counterexample to the TxGraffiti conjecture for every $k\ge1$. The graphs are the incidence graphs of minimally unsatisfiable $3$-CNF formulas in which every literal occurs exactly twice, a construction that appears in~\cite{Gupta}; the formulas are obtained by a gluing operation from a complete census of the eight such formulas of deficiency three on seven variables.
\end{abstract}

\section{Statement}
For a graph $G$, $\gamma(G)$ is the domination number, $i(G)$ the independent domination number (the minimum size of a maximal independent set), and $\gamma_e(G)=\mu^*(G)$ the edge domination number, the minimum size of a maximal matching. Since a maximal independent set is dominating, $\gamma(G)\le i(G)$. Two conjectures compare these to $\gamma_e$.
\begin{itemize}
\item[(A)] \cite{BFHMR}: if $G$ is $\Delta$-regular with $\Delta\ge1$, then $\gamma(G)\le\gamma_e(G)$. Known for $\Delta\le2$ trivially, for $\Delta\ge7$ by the Lov\'asz Local Lemma~\cite{Gupta}, for claw-free and fork-free graphs, and for every cubic graph on at most $48$ vertices~\cite[Prop.~7]{Gupta}; the cubic case is described as open in~\cite{Gupta}.
\item[(B)] \cite[Conj.~3]{DBP}: if $G$ is $r$-regular with $r>0$, then $i(G)\le\mu^*(G)$, and this bound is sharp. Generated by TxGraffiti in 2020, reported open for every $r\ge3$; (B) implies (A).
\end{itemize}
\begin{theorem}\label{main}
There are six pairwise non-isomorphic connected cubic graphs on $50$ vertices with $i(G)=16$ and $\mu^*(G)=15$; three of them have $\gamma(G)=16$. Hence (A) and (B) both fail for cubic graphs, and the three graphs are counterexamples to (A) of minimum order. Moreover, for every $r\equiv0\pmod3$ there is a connected $r$-regular graph with $i(G)>\mu^*(G)$.
\end{theorem}
The inequalities are attained with equality on many graphs ($K_{3,3}$, the Petersen graph, the prism over a triangle for (A), and the graph $G_{50}$ of Remark~\ref{gupta} for both), so only the inequalities fail, not the sharpness clauses. Conjecture (B) had been verified exhaustively for cubic graphs on at most $20$ vertices at minimum maximal matchings; the smallest cubic counterexample to (B) therefore has between $22$ and $50$ vertices. For (A) the minimum is exactly $50$.

\section{From formulas to graphs}
A CNF formula is \emph{minimally unsatisfiable} (MU) if it is unsatisfiable and every proper subset is satisfiable~\cite{KB}. Let $F$ be a $3$-uniform CNF formula on variables $x_1,\dots,x_v$ with $m$ clauses in which every literal occurs exactly twice (so $3m=4v$). Define $G(F)$ as follows: for each variable $x$ take two adjacent vertices $x^+,x^-$ (a \emph{pivot pair}); for each clause $C$ take a vertex $u_C$ adjacent to the endpoint named by each of its three literals ($x\mapsto x^+$, $\bar x\mapsto x^-$).
\begin{lemma}\label{dict}
$G(F)$ is a simple $3$-regular graph on $2v+m$ vertices; the $v$ pivot edges form a maximal matching $M$; and $F$ is unsatisfiable if and only if no independent set of pivot endpoints dominates all clause vertices. If $F$ is MU then $G(F)$ is connected.
\end{lemma}
\begin{proof}
Each endpoint has one pivot edge and two clause edges; each clause vertex has three edges to endpoints of three distinct pairs. The clause vertices are pairwise non-adjacent, so $M$ is maximal. An independent set of endpoints contains at most one endpoint per pair, i.e. it is a partial truth assignment, and it dominates $u_C$ iff it makes a literal of $C$ true. Connectedness follows since an MU formula has a connected clause--variable incidence graph.
\end{proof}
Thus for unsatisfiable $F$, no independent dominating set of $G(F)$ can be built by extending the clause vertices, and $|M|=v$ competes with $i(G(F))$ with no further structure to help; $\mu^*(G(F))\le v$ always.

\section{Gluing minimally unsatisfiable formulas}
\begin{lemma}[Gluing]\label{glue}
Let $A,B$ be MU formulas on disjoint variable sets, $a_1\ne a_2\in A$, $b_1\ne b_2\in B$, and $x$ a fresh variable. Then $K:=(A$ with $x$ added to $a_1,a_2)\cup(B$ with $\bar x$ added to $b_1,b_2)$ is MU. Literal occurrences are unchanged except that $x$ and $\bar x$ occur twice each, and $|K|-|\mathrm{var}\,K|=(|A|-|\mathrm{var}\,A|)+(|B|-|\mathrm{var}\,B|)-1$.
\end{lemma}
\begin{proof}
$K[x{=}1]\supseteq B$ and $K[x{=}0]\supseteq A$, so $K$ is unsatisfiable. Remove a clause of $A$'s side and set $x=0$: the $\bar x$-clauses are satisfied, $A$ minus that clause is satisfiable, and $B\setminus\{b_1,b_2\}$ is satisfiable as a proper subset of an MU formula; the variable sets are disjoint. The other side is symmetric.
\end{proof}
Let $R_3$ be the following formula on $7$ variables (two $2$-clauses and eight $3$-clauses, every literal exactly twice, MU):
\begin{align*}R_3=\;&(x_{1}\vee x_{2}) \; (x_{3}\vee x_{4}) \; (x_{1}\vee \bar x_{2}\vee x_{5}) \; (\bar x_{1}\vee x_{2}\vee x_{6}) \; (\bar x_{1}\vee x_{5}\vee x_{7})\\ &(\bar x_{2}\vee x_{6}\vee \bar x_{7}) \; (x_{3}\vee \bar x_{4}\vee \bar x_{5}) \; (\bar x_{3}\vee x_{4}\vee \bar x_{6}) \; (\bar x_{3}\vee \bar x_{5}\vee x_{7}) \; (\bar x_{4}\vee \bar x_{6}\vee \bar x_{7})\end{align*}
It is one of exactly eight isomorphism classes of MU formulas with these parameters, found by a complete canonical enumeration (Appendix~B). Gluing two copies of $R_3$ through a fresh variable $x_{15}$ placed in their four $2$-clauses gives
\begin{align*}K_1=\;&(x_{1}\vee x_{2}\vee x_{15})\quad (x_{3}\vee x_{4}\vee x_{15})\quad (x_{1}\vee \bar x_{2}\vee x_{5})\quad (\bar x_{1}\vee x_{2}\vee x_{6})\\
&(\bar x_{1}\vee x_{5}\vee x_{7})\quad (\bar x_{2}\vee x_{6}\vee \bar x_{7})\quad (x_{3}\vee \bar x_{4}\vee \bar x_{5})\quad (\bar x_{3}\vee x_{4}\vee \bar x_{6})\\
&(\bar x_{3}\vee \bar x_{5}\vee x_{7})\quad (\bar x_{4}\vee \bar x_{6}\vee \bar x_{7})\quad (x_{8}\vee x_{9}\vee \bar x_{15})\quad (x_{10}\vee x_{11}\vee \bar x_{15})\\
&(x_{8}\vee \bar x_{9}\vee x_{12})\quad (\bar x_{8}\vee x_{9}\vee x_{13})\quad (\bar x_{8}\vee x_{12}\vee x_{14})\quad (\bar x_{9}\vee x_{13}\vee \bar x_{14})\\
&(x_{10}\vee \bar x_{11}\vee \bar x_{12})\quad (\bar x_{10}\vee x_{11}\vee \bar x_{13})\quad (\bar x_{10}\vee \bar x_{12}\vee x_{14})\quad (\bar x_{11}\vee \bar x_{13}\vee \bar x_{14})\end{align*}
a $3$-uniform MU formula on $15$ variables with $20$ clauses in which every literal occurs exactly twice.

\section{The graph $G_1$ and its verification}
Let $G_1=G(K_1)$: $50$ vertices, $75$ edges, connected, cubic. Its adjacency list is in Appendix~A (vertices $2(j-1),2(j-1)+1$ are $x_j^{\pm}$; vertices $30+t$ are the clause vertices in the order above).
\begin{proposition}
$i(G_1)=16$, $\gamma(G_1)=14$ and $\mu^*(G_1)=15$.
\end{proposition}
\begin{proof}[Verification]
The $15$ pivot edges form a maximal matching (Lemma~\ref{dict}), so $\mu^*(G_1)\le15$. Conversely, in a cubic graph every edge of a maximal matching dominates at most five edges (itself and two at each endpoint), and a maximal matching dominates all edges, so any maximal matching of a cubic graph with $75$ edges has at least $\lceil75/5\rceil=15$ edges; hence $\mu^*(G_1)=15$ (also confirmed by integer programming). The set $\{1,3,4,7,9,11,12,15,17,19,20,23,24,27,30,40\}$ is independent and dominating, so $i(G_1)\le16$. That no independent dominating set of size $15$ exists was established by four independent methods: two SAT solvers (CaDiCaL~1.5.3 with a sequential-counter cardinality encoding; Glucose~4 with a totalizer encoding), the HiGHS integer programming solver, and a solver-free exhaustive branch-and-bound over independent dominating sets ($186{,}188$ search nodes). The domination number $14$ was computed by integer programming and SAT.
\end{proof}
The same construction applied to the other pairs of blocks gives, among the $36$ gluings of the eight blocks, exactly six graphs with $i=16$ (all with $\mu^*=15$), $25$ with $i=\mu^*=15$ and five with $i=14$. The six are pairwise non-isomorphic; their domination numbers are listed as well. For the three graphs with $\gamma=16$, the non-existence of a dominating set of size $15$ was verified by the HiGHS integer programming solver, by Glucose~4 with a totalizer encoding, and by a solver-free exhaustive branch-and-bound over dominating sets ($27{,}699{,}453$, $24{,}899{,}125$ and $25{,}088{,}473$ search nodes respectively); dominating sets of size $16$ were checked by definition. In all six graphs (and in $G_{50}$) the pivot matching is the unique maximal matching of size $15$, so it is the unique minimum maximal matching, and it admits no dominating transversal in the sense of~\cite{Gupta}. Every clause vertex is the centre of an induced claw and the graphs contain induced forks, so the claw-free and fork-free cases of (A), which are theorems, are not contradicted; the multiplicative bound $\gamma\le(7/6-1/204)\gamma_e$ of~\cite{BFHMR} is respected, $16/15<1.16$.
\begin{center}\begin{tabular}{lllll}\toprule graph & $i$ & $\gamma$ & $\mu^*$ & graph6\\\midrule
$R_3+R_3$ & 16 & 14 & 15 & \texttt{q`?G?C??G??@????\_???@?????G?????C??????G??????@????????\_???????@?????????L????OS???QO\_???WA???A@C???A@G???AS????@`?????aO????AI??????I?@????A\_C???@GO?????oC?????OG\_????@?c?????Cg??????KG??????PG??????CS????}\\
$R_3+R_6$ & 16 & 15 & 15 & \texttt{q`?G?C??G??@????\_???@?????G?????C??????G??????@????????\_???????@?????????L????OS???QO\_???WA???A@C???A@G???AS????@`?????aO????AI??????I?@????Q?C????p?????@CO?????Ca?????@@O?????AP??????GI??????DG??????Cc????}\\
$R_3+R_7$ & 16 & 15 & 15 & \texttt{q`?G?C??G??@????\_???@?????G?????C??????G??????@????????\_???????@?????????L????OS???QO\_???WA???A@C???A@G???AS????@`?????aO????AI??????I?@????Q?C????p?????@CO?????Ca?????@@G?????A`??????HG??????Cc??????@S????}\\
$R_6+R_6$ & 16 & 16 & 15 & \texttt{q`?G?C??G??@????\_???@?????G?????C??????G??????@????????\_???????@?????????L????QO???Pa????aG????cO???AA\_???@G\_???@@O????IO????AQ??????I?@????Q?C????p?????@CO?????Ca?????@@O?????AP??????GI??????DG??????Cc????}\\
$R_6+R_7$ & 16 & 16 & 15 & \texttt{q`?G?C??G??@????\_???@?????G?????C??????G??????@????????\_???????@?????????L????QO???Pa????aG????cO???AA\_???@G\_???@@O????IO????AQ??????I?@????Q?C????p?????@CO?????Ca?????@@G?????A`??????HG??????Cc??????@S????}\\
$R_7+R_7$ & 16 & 16 & 15 & \texttt{q`?G?C??G??@????\_???@?????G?????C??????G??????@????????\_???????@?????????L????QO???Pa????aG????cO???AAO???@O\_???@H?????HG?????i??????I?@????Q?C????p?????@CO?????Ca?????@@G?????A`??????HG??????Cc??????@S????}\\
\bottomrule\end{tabular}\end{center}

\section{Higher regularity}
For a graph $G$ and $k\ge1$ let $G[\overline K_k]$ be the lexicographic product with the edgeless graph on $k$ vertices: each vertex $v$ is replaced by $k$ pairwise non-adjacent copies, and copies of adjacent vertices are completely joined. If $G$ is connected and $r$-regular, $G[\overline K_k]$ is connected and $rk$-regular.
\begin{lemma}
$i(G[\overline K_k])=k\,i(G)$ and $\mu^*(G[\overline K_k])\le k\,\mu^*(G)$.
\end{lemma}
\begin{proof}
Let $I'$ be an independent dominating set of $G[\overline K_k]$ and $I$ the set of vertices of $G$ having a copy in $I'$. If a copy $(v,j)$ lies in $I'$ and another copy $(v,j')$ does not, then $(v,j')$ is dominated by some $(u,i)\in I'$ with $u\sim v$, but then $(u,i)\sim(v,j)$, contradicting independence; so $I'=I\times[k]$. Independence and domination of $I'$ transfer to $I$, hence $|I'|\ge k\,i(G)$, and $I\times[k]$ for an optimal $I$ shows equality. For the matching, take a maximal matching $M$ of $G$ and the $k$ edges $(u,j)(v,j)$, $j\in[k]$, for each $uv\in M$; the uncovered vertices are the copies of the $M$-uncovered vertices, an independent set, so this matching of size $k|M|$ is maximal.
\end{proof}
Hence $i(G_1[\overline K_k])=16k>15k\ge\mu^*(G_1[\overline K_k])$ for every $k$, which proves the last claim of Theorem~\ref{main}; for $k=2$ we verified $i=32$, $\mu^*=30$ directly by integer programming. No such extension exists for $\gamma$: a dominating set of $G[\overline K_k]$ may use a single copy of each vertex of a total dominating set of $G$, so $\gamma(G[\overline K_k])$ stays bounded while $\gamma_e$ grows, in accordance with the theorem of~\cite{Gupta} that (A) holds for all $\Delta\ge7$. Regularities $r\not\equiv0\pmod 3$ remain open: the natural analogue for $r=4$ is a $4$-uniform MU formula in which every literal occurs exactly three times (then $4m=6v$, i.e.\ $2m=3v$), whose realization is a $4$-regular graph with an obstructed maximal matching, and the smallest conceivable such formula has $12$ variables and $18$ clauses.

\section{Remarks}
\begin{remark}\label{gupta}
Unsatisfiable $3$-CNF formulas in which every literal occurs exactly twice are known to exist~\cite{BKS,BKS2}, and the least number of clauses of such a formula is $20$~\cite{ZPS}; it is this fact that yields the $48$-vertex bound of~\cite[Prop.~7]{Gupta}, and it shows that $n=50$ is the least order at which the present construction can produce anything, since $3m=4v$ and $m\ge20$ force $v\ge15$. Gupta~\cite[Thm.~2]{Gupta} builds the incidence graph of one such minimum formula and shows that its unique minimum maximal matching admits no dominating transversal while $\gamma=14<15=\gamma_e$. That graph is isomorphic to our $G_{50}=G(R_0\oplus R_0)$, the gluing of two copies of $R_0$ (itself the gluing of two copies of $F_3$), for which also $i=15=\mu^*$; the counterexamples arise from the other blocks of the census.
\end{remark}
\begin{remark}
In the language of minimally unsatisfiable formulas: every MU formula of deficiency one with literal degrees at most two has at least four more literal occurrences than clauses, every such formula of deficiency two has at least three, and no $3$-uniform MU formula with literal degrees at most two has deficiency $\le3$ (the last is a short descent argument using the pocket structure of restriction cores). The blocks $R_0,\dots,R_7$ are the complete list of MU formulas of deficiency three and defect two on seven variables; gluing produces $3$-uniform members at every deficiency $\ge5$, and by~\cite{ZPS} none exists at deficiency four with fewer than $20$ clauses.
\end{remark}
\begin{remark}
The route to these graphs went through a structural theory of obstructed maximal matchings in cubic graphs (demand systems, forcing, the deficiency laws above); the counterexamples live in the one anatomy---all pivot pairs, all loads two, no cycles---that the all-load-one theory and Haxell's independent transversal theorem~\cite{Haxell} exclude, which is why they were invisible to searches guided by that theory. Regularities $r\not\equiv0\pmod3$ remain open for (B), and degrees $4\le\Delta\le6$ for (A); the natural analogue for $r=4$ is a $4$-uniform MU formula in which every literal occurs exactly three times.
\end{remark}
\appendix
\section{Adjacency list of $G_1$}
\noindent\small
0:\ 1, 30, 32\\
1:\ 0, 33, 34\\
2:\ 3, 30, 33\\
3:\ 2, 32, 35\\
4:\ 5, 31, 36\\
5:\ 4, 37, 38\\
6:\ 7, 31, 37\\
7:\ 6, 36, 39\\
8:\ 9, 32, 34\\
9:\ 8, 36, 38\\
10:\ 11, 33, 35\\
11:\ 10, 37, 39\\
12:\ 13, 34, 38\\
13:\ 12, 35, 39\\
14:\ 15, 40, 42\\
15:\ 14, 43, 44\\
16:\ 17, 40, 43\\
17:\ 16, 42, 45\\
18:\ 19, 41, 46\\
19:\ 18, 47, 48\\
20:\ 21, 41, 47\\
21:\ 20, 46, 49\\
22:\ 23, 42, 44\\
23:\ 22, 46, 48\\
24:\ 25, 43, 45\\
25:\ 24, 47, 49\\
26:\ 27, 44, 48\\
27:\ 26, 45, 49\\
28:\ 29, 30, 31\\
29:\ 28, 40, 41\\
30:\ 0, 2, 28\\
31:\ 4, 6, 28\\
32:\ 0, 3, 8\\
33:\ 1, 2, 10\\
34:\ 1, 8, 12\\
35:\ 3, 10, 13\\
36:\ 4, 7, 9\\
37:\ 5, 6, 11\\
38:\ 5, 9, 12\\
39:\ 7, 11, 13\\
40:\ 14, 16, 29\\
41:\ 18, 20, 29\\
42:\ 14, 17, 22\\
43:\ 15, 16, 24\\
44:\ 15, 22, 26\\
45:\ 17, 24, 27\\
46:\ 18, 21, 23\\
47:\ 19, 20, 25\\
48:\ 19, 23, 26\\
49:\ 21, 25, 27\\
\normalsize
\section{The eight blocks}
Isomorphism classes of MU formulas on seven variables with two $2$-clauses, eight $3$-clauses and every literal occurring exactly twice ($R_0$ is the gluing of two copies of $F_3$; $R_4$ was found independently by a second engine):
\begin{description}\small
\item[$R_0$] $(x_{1}\vee x_{2}) \; (x_{3}\vee x_{4}) \; (x_{1}\vee \bar x_{2}\vee x_{5}) \; (\bar x_{1}\vee x_{2}\vee \bar x_{5}) \; (\bar x_{1}\vee x_{5}\vee x_{6}) \; (\bar x_{2}\vee \bar x_{5}\vee x_{6}) \; (x_{3}\vee \bar x_{4}\vee x_{7}) \; (\bar x_{3}\vee x_{4}\vee \bar x_{7}) \; (\bar x_{3}\vee \bar x_{6}\vee x_{7}) \; (\bar x_{4}\vee \bar x_{6}\vee \bar x_{7})$
\item[$R_1$] $(x_{1}\vee x_{2}) \; (x_{3}\vee x_{4}) \; (x_{1}\vee \bar x_{2}\vee x_{5}) \; (\bar x_{1}\vee x_{2}\vee x_{6}) \; (\bar x_{1}\vee \bar x_{3}\vee x_{5}) \; (\bar x_{2}\vee \bar x_{4}\vee x_{6}) \; (\bar x_{3}\vee \bar x_{5}\vee x_{7}) \; (x_{3}\vee \bar x_{6}\vee x_{7}) \; (x_{4}\vee \bar x_{5}\vee \bar x_{7}) \; (\bar x_{4}\vee \bar x_{6}\vee \bar x_{7})$
\item[$R_2$] $(x_{1}\vee x_{2}) \; (x_{3}\vee x_{4}) \; (x_{1}\vee \bar x_{2}\vee x_{5}) \; (\bar x_{1}\vee x_{2}\vee x_{6}) \; (\bar x_{1}\vee \bar x_{3}\vee x_{5}) \; (\bar x_{2}\vee x_{6}\vee x_{7}) \; (x_{3}\vee \bar x_{4}\vee \bar x_{7}) \; (\bar x_{3}\vee \bar x_{5}\vee \bar x_{7}) \; (x_{4}\vee \bar x_{5}\vee \bar x_{6}) \; (\bar x_{4}\vee \bar x_{6}\vee x_{7})$
\item[$R_3$] $(x_{1}\vee x_{2}) \; (x_{3}\vee x_{4}) \; (x_{1}\vee \bar x_{2}\vee x_{5}) \; (\bar x_{1}\vee x_{2}\vee x_{6}) \; (\bar x_{1}\vee x_{5}\vee x_{7}) \; (\bar x_{2}\vee x_{6}\vee \bar x_{7}) \; (x_{3}\vee \bar x_{4}\vee \bar x_{5}) \; (\bar x_{3}\vee x_{4}\vee \bar x_{6}) \; (\bar x_{3}\vee \bar x_{5}\vee x_{7}) \; (\bar x_{4}\vee \bar x_{6}\vee \bar x_{7})$
\item[$R_4$] $(x_{1}\vee x_{2}) \; (\bar x_{1}\vee x_{3}) \; (x_{1}\vee \bar x_{2}\vee x_{3}) \; (\bar x_{1}\vee x_{2}\vee x_{4}) \; (\bar x_{2}\vee \bar x_{3}\vee x_{4}) \; (\bar x_{3}\vee x_{5}\vee x_{6}) \; (\bar x_{4}\vee \bar x_{5}\vee x_{7}) \; (\bar x_{4}\vee \bar x_{6}\vee \bar x_{7}) \; (x_{5}\vee \bar x_{6}\vee x_{7}) \; (\bar x_{5}\vee x_{6}\vee \bar x_{7})$
\item[$R_5$] $(x_{1}\vee x_{2}) \; (\bar x_{1}\vee x_{3}) \; (x_{1}\vee \bar x_{2}\vee x_{3}) \; (\bar x_{1}\vee x_{2}\vee x_{4}) \; (\bar x_{2}\vee x_{4}\vee x_{5}) \; (\bar x_{3}\vee \bar x_{5}\vee x_{6}) \; (\bar x_{3}\vee \bar x_{6}\vee x_{7}) \; (\bar x_{4}\vee x_{5}\vee \bar x_{7}) \; (\bar x_{4}\vee x_{6}\vee x_{7}) \; (\bar x_{5}\vee \bar x_{6}\vee \bar x_{7})$
\item[$R_6$] $(x_{1}\vee x_{2}) \; (\bar x_{1}\vee x_{3}) \; (\bar x_{1}\vee x_{2}\vee x_{4}) \; (x_{1}\vee x_{3}\vee x_{5}) \; (\bar x_{2}\vee x_{4}\vee x_{6}) \; (\bar x_{2}\vee \bar x_{5}\vee \bar x_{6}) \; (\bar x_{3}\vee x_{5}\vee x_{7}) \; (\bar x_{3}\vee \bar x_{6}\vee \bar x_{7}) \; (\bar x_{4}\vee \bar x_{5}\vee x_{7}) \; (\bar x_{4}\vee x_{6}\vee \bar x_{7})$
\item[$R_7$] $(x_{1}\vee x_{2}) \; (\bar x_{1}\vee x_{3}) \; (\bar x_{1}\vee x_{2}\vee x_{4}) \; (x_{1}\vee x_{3}\vee x_{5}) \; (\bar x_{2}\vee x_{4}\vee x_{6}) \; (\bar x_{2}\vee \bar x_{5}\vee x_{7}) \; (\bar x_{3}\vee \bar x_{4}\vee x_{7}) \; (\bar x_{3}\vee x_{5}\vee \bar x_{6}) \; (\bar x_{4}\vee x_{6}\vee \bar x_{7}) \; (\bar x_{5}\vee \bar x_{6}\vee \bar x_{7})$
\end{description}
\begin{thebibliography}{9}
\bibitem{BFHMR} J.~Baste, M.~F\"urst, M.~A.~Henning, E.~Mohr, D.~Rautenbach, Domination versus edge domination, \emph{Discrete Appl. Math.} 285 (2020) 343--349; arXiv:1906.10420.
\bibitem{BKS} P.~Berman, M.~Karpinski, A.~D.~Scott, Approximation hardness and satisfiability of bounded occurrence instances of SAT, ECCC report TR03-022, 2003.
\bibitem{BKS2} P.~Berman, M.~Karpinski, A.~D.~Scott, Approximation hardness of short symmetric instances of MAX-3SAT, ECCC report TR03-049, 2003.
\bibitem{Gupta} C.~Gupta, Domination versus edge domination in regular graphs of degree at least seven, arXiv:2608.22498, 2026.
\bibitem{ZPS} T.~Zhang, T.~Peitl, S.~Szeider, Small unsatisfiable $k$-CNFs with bounded literal occurrence, in: \emph{27th International Conference on Theory and Applications of Satisfiability Testing (SAT 2024)}, LIPIcs vol.~305, Schloss Dagstuhl -- Leibniz-Zentrum f\"ur Informatik, 2024, pp.~31:1--31:22.
\bibitem{DBP} R.~Davila, B.~Brimkov, R.~Pepper, \emph{In Reverie Together: Ten Years of Mathematical Discovery with a Machine Collaborator}, arXiv:2507.17780 [cs.DM], 2025.
\bibitem{Haxell} P.~E.~Haxell, A note on vertex list colouring, \emph{Combin. Probab. Comput.} 10 (2001) 345--347.
\bibitem{KB} H.~Kleine B\"uning, O.~Kullmann, Minimal unsatisfiability and autarkies, in: \emph{Handbook of Satisfiability}, IOS Press, 2009, ch.~11.
\end{thebibliography}
\end{document}
